\documentclass[11pt]{article}
\usepackage[margin=1in]{geometry}
\usepackage{amsmath,amssymb}
\usepackage{booktabs}
\usepackage{hyperref}
\usepackage{listings}
\usepackage{xcolor}
\usepackage{caption}

\lstset{
  basicstyle=\ttfamily\small,
  breaklines=true,
  columns=fullflexible,
  backgroundcolor=\color{gray!8},
  frame=single,
  language=Python
}

\title{Variance-Guided Progressive Rollout Allocation for\\
Sequence-Level Policy Optimization: An Applied Integration\\
of VIGOR into GSPO for Free-Form Information Extraction}

\author{Applied report on the GSPO/VERL information-extraction pipeline\\ \texttt{NER\_Training}}
\date{August 2026}

\begin{document}
\maketitle

\begin{abstract}
We describe the design, implementation, and empirical validation of a variance-guided
progressive rollout allocation strategy integrated into a production Group Sequence
Policy Optimization (GSPO) reinforcement-learning pipeline for free-form named-entity
and event extraction. Standard GSPO with a fixed rollout count per prompt plateaus
after supervised warm-start and two rounds of LoRA-capacity escalation, holding at a
validation reward of $\approx 0.709$ for hundreds of steps regardless of learning rate
or adapter placement. We integrate VIGOR (Jiang et al., 2026), a progressive
rollout-allocation algorithm that concentrates sampling budget on high
current-policy-uncertainty prompts across refinement rounds within a single training
step, directly into VERL's actor-rollout-advantage pipeline — reusing VERL's existing
group-relative advantage estimator (which already supports variable per-prompt group
sizes) rather than reimplementing advantage computation, and fixing a latent
minibatch-sizing assumption in VERL's actor-update path that a fixed-rollout-count
training loop never exercises. The integration is validated first as an isolated,
adversarially-reviewed library against mock rollouts, then live against the real
training stack. In production training, the resulting policy breaks the fixed-rollout
plateau, reaching a validation reward of $0.718$ (vs.\ $0.709$ ceiling) and an
offline-evaluated composite extraction score of $0.738$ — the best result obtained in
this project — with the largest gains concentrated in the two weakest sub-skills
(event topic-code and action-code classification), exactly the metrics the allocator
was designed to target.
\end{abstract}

\section{Introduction}

Reinforcement learning from verifiable rewards (RLVR) for structured generation tasks
— here, emitting a JSON object of named entities and CAMEO-coded events from a news
document — typically samples a fixed number of rollouts $G$ per prompt and computes a
group-relative advantage. This is the GRPO/GSPO family's default, and it is what our
production pipeline used through two escalation attempts (Sections~\ref{sec:baseline}).
Fixed-$G$ sampling spends equal budget on every prompt regardless of how much signal
that prompt currently offers: a prompt the policy already solves, or one it is
uniformly hopeless at, yields near-zero reward variance and therefore near-zero
gradient signal (Jiang et al., 2026, Section 2), while a prompt straddling the
decision boundary yields the largest usable signal. VIGOR reallocates budget toward
the latter by refining in rounds: start every active prompt with a small rollout
count, measure cumulative reward variance, keep only the highest-variance fraction,
double their allocation, and repeat.

This report documents an applied integration of that idea into a real,
already-in-production GSPO/VERL training loop for information extraction, not a
reimplementation of the RL framework. Section~\ref{sec:background} summarizes GSPO
and VIGOR. Section~\ref{sec:method} describes the integration architecture — in
particular, a simplification that removed an entire component of the original
algorithm from the applied system, and a defect in the host framework's actor-update
path that the integration surfaced and fixed. Section~\ref{sec:validation} describes
a two-stage validation methodology (isolated adversarial review, then live smoke
test) motivated by a real production incident earlier in the same pipeline in which
an unhandled type error in the reward function silently killed a 38-step run.
Section~\ref{sec:results} reports the production training outcome against two
fixed-rollout baselines. Section~\ref{sec:discussion} discusses limitations and an
unrelated but consequential operational failure (disk exhaustion from
unbounded checkpoint retention) encountered during the run.

\section{Background}
\label{sec:background}

\subsection{GSPO}

Group Sequence Policy Optimization (Zheng et al., 2025, arXiv:2507.18071) replaces
GRPO's per-token importance-ratio clipping with a single, sequence-level ratio and
clip decision per generated response. For a document $x$, $G$ completions
$y_1,\dots,y_G$ are sampled; rewards $r_i$ are scored by a task-specific verifier and
converted to group-relative advantages
\begin{equation}
A_i = \frac{r_i - \bar{r}}{\sigma_r + \delta}, \qquad
\bar{r} = \frac1G\sum_i r_i, \qquad
\sigma_r = \sqrt{\frac1G\sum_i (r_i-\bar r)^2}.
\end{equation}
The policy loss uses a sequence-level ratio $\rho_i(\theta)$ (rather than a per-token
ratio) with a single clip per sequence, aggregated as \texttt{seq-mean-token-mean}.
In our host framework (VERL), this corresponds to
\texttt{algorithm.adv\_estimator=grpo} (group-relative advantages, computed by
\texttt{core\_algos.compute\_grpo\_outcome\_advantage}) combined with
\texttt{actor\_rollout\_ref.actor.policy\_loss.loss\_mode=gspo} (the sequence-level
ratio/clip, in \texttt{compute\_policy\_loss\_gspo}). Both components are
production VERL code and are not modified by this work; the object of integration
is exclusively \emph{how many rollouts are generated for which prompts before}
these two functions are called.

\subsection{VIGOR}

VIGOR (Jiang, Liu, and Mirzasoleiman, 2026, arXiv:2607.22002) proves that, under
binary RLVR rewards with the KL term removed, the gradient-norm bound of the
GRPO-family objective is proportional to the within-group reward standard deviation
$\sigma$ — maximized when the policy is genuinely uncertain about a prompt (success
probability $p\approx 0.5$), and near-zero when a prompt is already solved or
hopeless. VIGOR's Algorithm 1 exploits this directly: starting from an active set
$A_1 = $ the full sampling batch, at each of $T$ refinement rounds it (a) generates
$m^{(t)}$ \emph{new} rollouts per currently-active prompt, (b) merges them into that
prompt's cumulative rollout set (never discarding earlier rounds), (c) computes the
population variance of each active prompt's cumulative rewards, and, if
$t<T$, (d) retains only the top-$\alpha$ fraction of the active set by variance and
multiplies the per-prompt rollout count by $\gamma$ for the survivors. With the
paper's own budget-conserving configuration $\alpha\gamma = 1$ (e.g.\ $\alpha=0.5,
\gamma=2$), total expected rollout consumption collapses to $C_{\text{total}} =
B\cdot m_0\cdot T$ — matching a fixed-rollout baseline of $G=m_0 T$ rollouts/prompt
at the same nominal cost, while spending that budget on a small, sharply skewed
subset of prompts rather than uniformly. Advantages for the final GRPO update are
computed \emph{once} per surviving prompt, over its entire cumulative rollout set,
not per refinement round — the single most consequential implementation detail
called out in the paper and independently re-derived in Section~\ref{sec:validation}.

\section{Method}
\label{sec:method}

\subsection{A key simplification: reusing the host framework's existing advantage estimator}
\label{sec:simplification}

A na\"ive port of VIGOR would reimplement its own group-relative advantage
computation, since Algorithm~1 as stated operates on an explicit
rollout-set-per-prompt data structure. We instead inspected VERL's own
group-advantage implementation and found it already satisfies VIGOR's requirements
without modification:

\begin{lstlisting}[caption={VERL's grpo advantage estimator, unmodified}, label=lst:grpo]
# core_algos.py (unmodified VERL source)
id2score = defaultdict(list)
for i in range(bsz):
    id2score[index[i]].append(scores[i])     # index[i] is a per-row "uid" column
# ... group-relative normalization per uid, with NO assumption
# that every uid appears the same number of times.
\end{lstlisting}

Grouping is keyed purely on a non-tensor \texttt{uid} column carried per row, with a
\texttt{defaultdict(list)} accumulation pattern that places no constraint on how many
rows share a \texttt{uid}. This is precisely VIGOR's variable-group-size output shape
(a solved prompt might retain 2 rollouts; a maximally-uncertain survivor of every
round retains $m_0(1+\gamma+\gamma^2+\dots) = 30$ under the paper's defaults). The
practical consequence: \textbf{the applied integration does not port VIGOR's
advantage-computation code at all.} It is responsible only for (i) deciding, each
round, which prompts continue and how many new rollouts they receive, and (ii)
assigning and preserving a consistent \texttt{uid} per prompt across rounds so that
step~(iii), VERL's own unmodified \texttt{compute\_advantage}, groups correctly. This
follows the architectural constraint articulated in the source algorithm's own
design notes: the allocator decides \emph{which prompts, how many rollouts, when to
stop generating}; the host trainer remains solely responsible for log-probabilities,
advantages, clipping, and the optimizer step.

\subsection{Integration architecture}

The production entry point (\texttt{grpo/run\_gspo\_vigor.py}) mirrors the host
framework's own startup sequence — instantiating the same trainer class
(\texttt{RayPPOTrainer}) and calling its own \texttt{init\_workers()} — and replaces
only the fixed-rollout \texttt{fit()} step body with a custom loop implementing
Listing~\ref{alg:vigor-step}.

\begin{lstlisting}[caption={VIGOR training step, as integrated}, label=alg:vigor-step]
assign uid per prompt in the training batch
active = all prompts;  collected = {};  m = m0;  budget_used = 0
for round in range(T):
    requested = len(active) * m
    if budget_used + requested > budget_target:
        shrink m first; if m < 1, truncate active (drop lowest-variance)
    new_rollouts = async_rollout_manager.generate_sequences(active, m)
    # policy snapshot is frozen: no weight-sync call occurs inside this loop
    score(new_rollouts) via reward_ie.compute_score   # or reuse agent-loop rm_scores
    merge new_rollouts into collected[prompt] for prompt in active
    budget_used += len(new_rollouts)
    if round < T - 1:
        utility[p] = population_variance(collected[p].rewards) for p in active
        keep = ceil(alpha * len(active))
        active = top_k(active, keep, key=utility, tie_break=uid)  # deterministic
        m = round(m * gamma)
batch = DataProto.concat(collected.values())   # ragged uid-group sizes, by design
batch = trainer._compute_old_log_prob(batch)
batch = compute_advantage(batch, adv_estimator="grpo")   # unmodified VERL call
trainer._update_actor(batch)                              # see Section 3.3
\end{lstlisting}

The policy-freeze invariant central to VIGOR's correctness — every round within a
step must sample from the \emph{same} policy, since refinement is meaningless if
the target is moving — was verified structurally rather than assumed: the host
framework connects actor weights to rollout replicas through exactly one call
(\texttt{checkpoint\_manager.update\_weights()}), which in the unmodified fixed-rollout
loop appears only once per step, after the actor update. The integrated loop
preserves this exactly — the call is never issued inside the round loop.

\subsection{A latent host-framework assumption this integration surfaced}
\label{sec:minibatch-fix}

VERL's actor-update path computes an effective PPO minibatch size as
\texttt{ppo\_mini\_batch\_size} $\times$ \texttt{rollout.n} — a quantity that is
exact under fixed-rollout sampling (every prompt contributes exactly
\texttt{rollout.n} rows) and is asserted, in two places, to divide the actual batch
row count exactly (once in the data-loading iterator, once in the distributed
training-engine dispatcher). It is additionally used, unmodified, as the loss
aggregation's normalizing denominator. A VIGOR batch's true row count is
$\sum_p |Y_p|$ over whatever prompts survived to whatever round — a number with no
fixed relationship to \texttt{rollout.n} — so this host assumption fails two
hard-coded divisibility assertions and, more subtly, would silently mis-scale the
gradient if the assertions were bypassed rather than fixed. We resolve this at the
integration boundary rather than by patching the host framework: before calling the
unmodified \texttt{\_update\_actor}, we compute the nearest valid divisor of the
\emph{actual} row count (a multiple of the data-parallel world size, floored to
avoid degenerating to minibatch size 1 on a prime-ish row count) and temporarily
override \texttt{rollout.n=1} together with the chosen minibatch size for the
duration of that one call, so the denominator VERL itself recomputes matches the
data it is actually operating on. We separately re-derived that this cannot corrupt
the GSPO sequence-level ratio/clip itself: that computation is entirely per-row and
never reads \texttt{uid} or any other cross-row grouping signal; the only place
group membership matters is the advantage computation of Section~\ref{sec:simplification},
which completes \emph{before} minibatching and is therefore unaffected by how the
minibatches are subsequently split.

\subsection{Reward attribution and the pilot pool}

Each round's rollouts are scored by the same verifier used for training
(\texttt{grpo/reward\_ie.py::compute\_score}), by default reusing the reward already
computed by the host framework's own agent-execution loop for the training reward
itself (\texttt{score\_source=auto}), so the variance signal driving selection is
bit-identical to the signal that will drive the eventual policy gradient — a
stronger guarantee than recomputing it independently would give. An alternate mode
recomputes the score directly from decoded text on the driver process for
configurations where the reward is not already attached to the rollout output; an
early version of this path read ground truth from the wrong data object (the
rollout-output object, which the host agent loop does not echo prompt-side fields
onto) and was corrected to read from the prompt-aligned object instead, with hard
failures replacing what had been a silent \texttt{None}.

\section{Validation methodology}
\label{sec:validation}

Two properties of this codebase's history motivated a specifically two-stage
validation process rather than either static review alone or a direct production
run: (1) an earlier, unrelated defect in the reward function — a truthy non-iterable
model output silently escaping an \texttt{x or []} guard — killed a full production
run via an unhandled exception propagating through the distributed reward worker,
demonstrating that plausible-looking reward-adjacent code can hide crash-on-first-
malformed-input defects that pass casual review; and (2) rollout-allocation logic of
this kind (cumulative merging across rounds, variance-based selection, budget
accounting) is exactly the class of logic the source algorithm's own design notes
flag as prone to \emph{silent} correctness failures — wrong output, not a crash —
under a policy-freeze violation, an off-by-one in cross-round merging, or advantage
normalization computed per-round instead of over the cumulative group.

\textbf{Stage 1 — isolated adversarial validation.} Before any integration code was
written, VIGOR (and, separately, Pilot-Commit) were implemented as a standalone
library against mock rollout generation, with an independent reviewer explicitly
tasked to \emph{refute} each of seven correctness properties (policy freeze, uid/
group integrity, cumulative-not-discarded merging, current-active-set-relative
selection, budget conservation, advantage normalization over the full cumulative
group, and — specific to Pilot-Commit — replay-buffer/staleness handling) by
independently re-deriving the relevant code paths rather than trusting the
implementer's self-report. This surfaced zero defects in VIGOR and six real defects
in Pilot-Commit (Section~\ref{sec:pc-bugs}), establishing VIGOR as the trustworthy
candidate for integration and motivating deferring Pilot-Commit.

\textbf{Stage 2 — live smoke test against the real training stack.} The same
adversarial process was then applied to the \emph{integration} code itself
(Section~\ref{sec:method}), surfacing four further defects (Section~\ref{sec:bugs}),
followed by a mandatory live execution — a handful of training steps against a tiny
dataset and the production checkpoint, GPU-resident, not mocked — on the explicit
premise that static review, however thorough, cannot substitute for observed
execution of training-loop code. The first three live attempts failed for reasons
unrelated to the code under test (an external, unrelated GPU tenant left
insufficient device memory); the fourth succeeded, producing round-by-round
diagnostics that matched the theoretical budget-conserving schedule exactly.

\subsection{Defects found in the standalone Pilot-Commit implementation}
\label{sec:pc-bugs}

For completeness, since Pilot-Commit remains unintegrated but its standalone
correctness bears on future work: a replay-buffered prompt was not excluded from
the next pilot pool, so it was silently re-piloted and the resulting fresh state
was overwritten by the stale buffered one — discarding paid-for on-policy rollouts
in favor of stale ones; a prompt selected from a fresh pilot could retain a stale
buffer entry and be trained on twice, once with outdated rollouts; the staleness
clock (\texttt{max\_policy\_delay}) never advanced under the base call signature,
silently disabling the staleness guard entirely; the default survivor-ranking policy
re-imposed a preference for success-probability near $0.5$ that the source paper
explicitly argues against; and one ranking policy was a no-op due to a
tie-in-sort-key artifact. All six were fixed and covered by regression tests
targeting the exact adversarial probes that found them.

\subsection{Defects found in the production integration}
\label{sec:bugs}

Reward-write ordering: a round's score was attached to the rollout-output object
\emph{after} that object had already been merged (via tensor-dictionary union) into
the batch that would later be searched for it — the merge operation copies existing
keys at the time it runs, so a key written afterward is invisible to the merged
object, and the very first training step would raise a lookup error. Fixed by
scoring before merging. Ground-truth source: the alternate direct-scoring path
(Section~\ref{sec:method}) read prompt metadata from the rollout-output object
rather than the prompt-aligned input object; fixed with hard failure guards
replacing a previous silent \texttt{None}. A diagnostic warning tested a
configuration node that is a real, always-populated part of the composed
configuration (hence never \texttt{None}) rather than the specific sub-switches
that actually enable the feature it warns about, causing it to fire on every run
regardless of configuration; corrected to test the real switches. A documented
example command in the operator-facing wrapper script was itself non-executable
against the host framework's own configuration validator; corrected.

\section{Experimental setup}

\textbf{Base model and adaptation.} Qwen3.5-0.8B, text-only checkpoint (see companion
process document for the architecture-specific extraction steps this required),
LoRA rank 16 on the two full-attention decoder layers (indices 19 and 23 of 24;
Qwen3.5 interleaves 18 Gated-DeltaNet linear-attention layers with 6 standard
full-attention layers), initialized from a checkpoint that had already completed
supervised warm-start and one prior round of fixed-rollout GSPO training on this
same adapter configuration — i.e.\ the comparison in Section~\ref{sec:results} is
\emph{continued} training under the two allocation strategies from an identical
starting point, not two runs from a common but earlier ancestor.

\textbf{VIGOR configuration.} $m_0=2$, $T=4$, $\gamma=2.0$, $\alpha=0.5$ (the source
paper's own reported configuration; $\alpha\gamma=1$, the budget-conserving
regime). Per-step budget $= B\cdot m_0\cdot T = 16\times2\times4=128$ rollout rows,
against a fixed-rollout baseline of $16\times4=64$ rows/step — i.e.\ VIGOR is
compared here at \emph{twice} the fixed baseline's nominal per-step rollout cost, a
deliberate choice consistent with the source paper's own experimental protocol of
comparing allocation strategies at matched \emph{cumulative}, not per-step, budget
over a full run.

\textbf{Optimizer and RL hyperparameters.} Identical across both conditions:
\texttt{lr}$=1\times10^{-6}$, GSPO clip range $[3\times10^{-4}, 4\times10^{-4}]$,
\texttt{loss\_agg\_mode=seq-mean-token-mean}, LoRA rank 16, no KL penalty.

\textbf{Evaluation.} Validation reward reported by the training loop's own
periodic validation pass against a held-out 196-document sample every 25 steps
(the aggregate weighted score of Section~1's reward function). A separate,
higher-fidelity offline evaluation (15 held-out documents, full sub-metric
breakdown, human-inspected raw JSON output) was run at selected checkpoints to
verify the aggregate number reflects genuine extraction quality rather than
reward-function exploitation — motivated directly by the Section~\ref{sec:results}
finding that an earlier checkpoint's climbing aggregate reward had in fact
reflected exactly such exploitation (Section 3 of the companion process document).

\section{Results}
\label{sec:results}

\subsection{Validation reward over training}

\begin{table}[h]
\centering
\caption{Validation reward, fixed-rollout GSPO baseline vs.\ VIGOR, from an
identical starting checkpoint (val reward $0.7090$).}
\label{tab:val-trend}
\begin{tabular}{@{}rrr@{}}
\toprule
Step & Fixed-rollout baseline (best of two escalation tiers) & VIGOR \\
\midrule
0    & 0.7076 -- 0.7090 & 0.7064 \\
75   & 0.7050 -- 0.7081 & 0.7066 \\
150  & 0.7063 -- 0.7092 & 0.7092 \\
250  & --                & 0.7129 \\
400  & --                & \textbf{0.7158} \\
875  & --                & \textbf{0.7181 (peak)} \\
950  & --                & 0.7172 (best \emph{saved} checkpoint) \\
\bottomrule
\end{tabular}
\end{table}

The fixed-rollout baseline, across two LoRA-capacity escalation tiers and one
learning-rate probe, never exceeded a validation reward of $0.7092$ in over 400
combined training steps, oscillating in a tight $0.705$--$0.709$ band. VIGOR
matched that ceiling by step 150 and continued climbing — non-monotonically, as
expected of a noisy 196-document validation sample, but with a clear upward trend
— reaching $0.7181$ at step 875, a sustained $\sim$1-percentage-point absolute
improvement over the fixed-rollout ceiling that held for hundreds of subsequent
steps rather than reverting.

Notably, at matched \emph{step count} in the early regime ($\le$75 steps), VIGOR
underperformed the fixed-rollout baseline slightly (e.g.\ offline composite score
$0.702$ vs.\ $0.730$ at step 70) despite consuming twice the rollout budget per
step — consistent with the source paper's framing that the algorithm's benefit is
a \emph{cumulative-budget-efficiency} claim over a full run, not a per-step
comparison, and we report this contemporaneous negative-looking result rather than
omitting it.

\subsection{Offline composite evaluation}

\begin{table}[h]
\centering
\caption{Sub-metric breakdown, offline 15-document evaluation, best checkpoints
under each condition.}
\label{tab:submetrics}
\begin{tabular}{@{}lrrrrr@{}}
\toprule
Checkpoint & score & entity F1 & subtype & topic F1 & action F1 \\
\midrule
SFT cold-start (no RL) & 0.701 & 0.774 & 0.911 & 0.360 & 0.60 \\
Fixed-rollout, best tier & 0.730 & 0.768 & 0.871 & 0.463 & 0.70 \\
VIGOR, step 1010 & \textbf{0.738} & 0.751 & 0.839 & \textbf{0.493} & \textbf{0.80} \\
\bottomrule
\end{tabular}
\end{table}

The VIGOR checkpoint attains the best composite score of any checkpoint evaluated
in this project, driven almost entirely by the two sub-metrics with the largest
remaining headroom (topic-code F1 and action-code F1) — precisely the components a
variance-guided allocator would be expected to improve, since a fixed-rollout
policy spends equal sampling budget regardless of a given document's topic/action
ambiguity, while VIGOR concentrates budget on exactly the documents where the
policy's current topic/action predictions are least certain. Entity F1 and subtype
accuracy are marginally lower than the fixed-rollout tier's best snapshot
($0.751$ vs.\ $0.768$; $0.839$ vs.\ $0.871$) — both already comfortably high and,
we hypothesize, receiving comparatively less marginal benefit from variance-guided
reallocation since entity extraction shows lower per-document reward variance to
exploit in the first place. All fifteen evaluated outputs used correct JSON
structure, correct field names, and real CAMEO-format topic/action codes, with
zero truncation or repetition failures — the schema-fidelity problem that
motivated the supervised warm-start (companion process document, Section 3) did
not recur under either RL condition once fixed.

\section{Discussion}
\label{sec:discussion}

\subsection{Why the host framework's advantage estimator was reusable}

The generality that made Section~\ref{sec:simplification}'s simplification possible
is not incidental: any RL framework implementing group-relative advantages by
grouping observed rollouts on an explicit group identifier, rather than by
positional or fixed-stride assumptions over the batch, is compatible with
variable-group-size allocators like VIGOR with zero modification to that
component. The corollary risk is exactly what Section~\ref{sec:minibatch-fix}
found: any \emph{other} component of the same framework that derives a
group-related quantity from the fixed-rollout-count assumption instead — here, a
minibatch-size product used as both a divisibility precondition and a loss-scaling
denominator — will fail, and may fail silently rather than loudly, under variable
group sizes. We recommend that any future integration of a variable-group-size
allocator into an unfamiliar RL framework specifically audit every downstream use
of the rollout-count hyperparameter, not only the advantage-computation call site,
before assuming compatibility.

\subsection{An operational failure orthogonal to the algorithm}

During the production run reported in Section~\ref{sec:results}, training halted
around step 1019 with no exception trace. Diagnosis attributed this to disk
exhaustion: per-step checkpointing at a 10-step interval had, over the course of
the run, accumulated over one hundred $\sim$4~GB checkpoints (385~GB) on a shared
volume, and a checkpoint write that coincided with the volume reaching capacity
left a truncated, unloadable checkpoint at the exact failure timestamp. This is
unrelated to the allocation algorithm itself but is reported here because it is a
direct consequence of a long-running RL job's storage growth rate not having been
considered a first-class parameter alongside the RL hyperparameters — a
generalizable lesson for any multi-day RL training run with frequent
checkpointing, independent of which allocation strategy is used.

\subsection{Limitations}

The reported validation-reward comparison (Table~\ref{tab:val-trend}) uses a
196-document held-out sample and is visibly noisy step-to-step; the reported trend
is directional and sustained rather than a single clean monotonic curve, and a
larger held-out set would tighten the comparison. The offline composite evaluation
(Table~\ref{tab:submetrics}) uses only 15 documents per checkpoint, chosen for
evaluation speed on CPU-only inference to avoid contending with concurrent GPU
training; absolute sub-metric values at this sample size carry meaningful sampling
variance, though the qualitative pattern (schema fidelity uniformly correct;
topic/action codes the primary beneficiary) was corroborated across multiple
checkpoints and sample draws over the course of the project. Pilot-Commit, the
companion algorithm to VIGOR in its source literature, was validated in isolation
but not integrated into production training; its combination with VIGOR (an
explicitly hybrid, non-paper-sanctioned strategy in the source literature) was
out of scope for this work.

\section{Conclusion}

We integrated a variance-guided progressive rollout-allocation algorithm into a
production sequence-level policy-optimization training loop for structured
information extraction, by reusing the host framework's existing group-relative
advantage machinery rather than reimplementing it, fixing a latent
fixed-rollout-count assumption in an unrelated part of the same framework that the
integration exposed, and validating the result through isolated adversarial review
followed by mandatory live execution rather than static review alone. In
production training, the resulting policy broke a validation-reward plateau that
two prior escalation strategies (supervised warm-start and LoRA-capacity expansion)
had failed to move, with the largest gains concentrated in exactly the sub-skills
the allocator was designed to prioritize.

\section*{References}
\begin{itemize}
\item Zheng, C., et al. \emph{Group Sequence Policy Optimization.} arXiv:2507.18071, 2025.
\item Jiang, H., Liu, H., Mirzasoleiman, B. \emph{Learning as Reasoning Unfolds:
  Progressive Rollout Allocation for Efficient Reinforcement Learning.}
  arXiv:2607.22002, 2026.
\item Kim, W., Yang, Z., Yan, J.N., Liu, J. \emph{Spend Your Rollouts Where It
  Counts: Rollout Allocation for Group-Based RL Post-Training.} arXiv:2605.26606, 2026.
\end{itemize}

\end{document}
